For an obstacle prescribed on a compact set, the same perturbation removes pluripolar exceptions without changing the regularized envelope.

\begin{proposition}\label{prop:compact-obstacle-polar-invariance}
Let $\theta$ be big, let $K\subseteq X$ be compact and non-pluripolar, and let $\psi$ be a bounded real function on $K$. Set
\[
 \begin{aligned}
 R&=\sup\{u\in\PSH(X,\theta):u\leq\psi\text{ on }K\},\\
 Q&=\sup\{u\in\PSH(X,\theta):u\leq\psi\text{ q.e. on }K\}.
 \end{aligned}
\]
Then $R^*=Q^*$ is a genuine $\theta$-plurisubharmonic function with minimal singularities. It is bounded if $\theta\geq0$. Removing any pluripolar subset of $K$ from the pointwise obstacle constraint also leaves this regularization unchanged.
\end{proposition}
\begin{proof}
Put $m=\inf_K\psi$ and $M=\sup_K\psi$. Since $V_\theta\leq0$, the function $V_\theta+m$ belongs to the pointwise family. Applying \cref{lem:compact-obstacle-upper-bound} with the constant obstacle $M$ gives a common upper bound $C$ for the larger family. The exact-background upper-envelope argument therefore makes $R^*$ and $Q^*$ genuine $\theta$-plurisubharmonic functions.

Fix a competitor $u$ for $Q$ and define a bounded function on $X$ by
\[
 f(x)=
 \begin{cases}
 \psi(x),&x\in K,\\
 \max\{\sup_Xu,M\},&x\notin K.
 \end{cases}
\]
Then $u\leq f$ quasi-everywhere. By \cref{lem:big-pointwise-competitor}, there are $u_t\in\PSH(X,\theta)$ satisfying $u_t\leq f$ everywhere and $u_t\to u$ almost everywhere. Each $u_t$ is a competitor for $R$, so $R^*\geq u$ almost everywhere. \Cref{lem:qpsh-almost-everywhere-order} gives this inequality everywhere. Taking the supremum and regularization yields $R^*\geq Q^*$; inclusion of the competitor families gives the reverse inequality.

Write $F=R^*=Q^*$. Since $F\leq C$, the potential $F-C$ is a competitor for $V_\theta$. Thus
\[
 V_\theta+m\leq F\leq V_\theta+C,
\]
which proves minimal singularities and, when $\theta\geq0$, boundedness.

If $E\subseteq K$ is pluripolar, the family constrained on $K\setminus E$ lies between the pointwise and quasi-everywhere families. Its regularization consequently equals $F$. This uses the upper bound on the original compact $K$, even when $K\setminus E$ is not compact.
\end{proof}
\index{Envelope!pluripolar obstacle exceptions}
